% ============================================================================
% SECTION 4: DUAL STRUCTURE AND BELIEF GEOMETRY
% ============================================================================
\section{Dual Structure, Boundary Limits, and Belief Geometry}\label{sec:dual}

\subsection{Regular dual flatness of the interior family}\label{subsec:dual-flat}

\begin{theorem}[Regular dual flatness]\label{thm:dual-flat-walsh}
The Walsh exponential family~\eqref{eq:walsh-exp} is a regular minimal
exponential family on $\mathcal X$.  Its interior carries the standard dually
flat information-geometric structure: natural coordinates $\theta$ are
e-affine, expectation coordinates $\eta$ are m-affine, the log-partition
function $\psi$ is the e-potential, and its Legendre dual is
$\varphi(\eta)=-H(p_\eta)$.
\end{theorem}

\begin{proof}
The natural parameter space is the open set $\R^{2^{2s}-1}$.  The sufficient
statistics are bounded, so the partition function is finite and smooth for all
$\theta$.  Minimality follows from the orthogonality and linear independence of
the nonconstant Walsh characters together with the constant character.  The
standard theorem for regular minimal exponential families then gives dual
flatness and the Legendre correspondence~\cite{amari2000methods,ay2017information}.
\end{proof}

\begin{remark}[Boundary discipline]
Theorem~\ref{thm:dual-flat-walsh} applies to full-support distributions only.
The exact key distributions $P_k$ lie in the closure.  Statements about
geodesics, Christoffel symbols, Bregman identities, or curvature at $P_k$ are
therefore meaningful only after regularization or after being rephrased as
finite boundary covariance statements.
\end{remark}

\subsection{KL divergence between exact key distributions}\label{subsec:kl-keys}

\begin{proposition}[Support separation]\label{prop:kl-keys}
If $S$ is bijective and $k\neq k'$, then
\begin{equation}\label{eq:kl-infinite}
  \operatorname{supp}(P_k)\cap\operatorname{supp}(P_{k'})=\varnothing,
  \qquad
  \KL(P_k\|P_{k'})=+\infty.
\end{equation}
\end{proposition}

\begin{proof}
If $(a,b)$ belongs to both supports, then
$S(a\oplus k)=b=S(a\oplus k')$.  Since $S$ is injective,
$a\oplus k=a\oplus k'$, hence $k=k'$, a contradiction.  Thus $P_{k'}$ assigns
zero probability to every point in the support of $P_k$ when $k\neq k'$, which
implies infinite KL divergence.
\end{proof}

\begin{remark}[No finite Bregman identity at singular endpoints]
The usual equality between KL divergence and the Bregman divergence of the dual
potential is an interior statement.  For $P_k$ and $P_{k'}$ with disjoint
supports, one may study
$\KL(P_k^\varepsilon\|P_{k'}^\varepsilon)$ and then let
$\varepsilon\downarrow0$, but the limiting divergence is infinite.  We therefore
do not use finite Pythagorean decompositions directly between distinct exact
key distributions.
\end{remark}

\subsection{Mixtures induced by key beliefs}\label{subsec:belief-mixture}

Let $Q=\Delta_{2^s-1}^{\circ}$ be the interior of the simplex over keys.  A
belief $q\in Q$ induces a distribution on input-output pairs
\begin{equation}\label{eq:mixture-belief}
  \bar P_q(a,b)=\sum_{k\in\B{s}}q(k)P_k(a,b)
  =2^{-s}\sum_{k\in\B{s}}q(k)\mathbf 1[b=S(a\oplus k)].
\end{equation}
If $q$ has full support, then $\bar P_q$ has full support on $\mathcal X$.

\begin{proposition}[Walsh representation of belief mixtures]\label{prop:belief-walsh}
Let
\[
  \widehat q(\alpha)=\sum_{k\in\B{s}}q(k)(-1)^{\alpha\cdot k}.
\]
Then
\begin{equation}\label{eq:belief-walsh}
  \bar P_q(a,b)
  =2^{-2s}\sum_{\alpha,\beta\in\B{s}}
  \widehat q(\alpha)\widehat S(\alpha,\beta)
  (-1)^{\alpha\cdot a\oplus\beta\cdot b}.
\end{equation}
In particular, the uniform belief induces the uniform distribution on
$\mathcal X$.
\end{proposition}

\begin{proof}
Apply Walsh inversion to the indicator
$\mathbf 1[b=S(a\oplus k)]$ and sum over $k$ with weights $q(k)$.  For uniform
$q$, $\widehat q(\alpha)=\mathbf 1[\alpha=0]$, and
$\widehat S(0,0)=2^s$ while $\widehat S(0,\beta)=0$ for $\beta\neq0$ because
$S$ is bijective.
\end{proof}

\subsection{Fisher--Rao geometry of the belief simplex}\label{subsec:belief-geometry}

\begin{definition}[Belief simplex metric]\label{def:belief-simplex}
For $q\in Q$ and tangent vectors $u,v$ with $\sum_k u_k=\sum_k v_k=0$, the
Fisher--Rao metric is
\begin{equation}\label{eq:belief-fisher}
  g_q^{(Q)}(u,v)=\sum_{k\in\B{s}}\frac{u_kv_k}{q(k)}.
\end{equation}
\end{definition}

Under the square-root map $q\mapsto(\sqrt{q(k)})_k$, this metric is the round
metric on a sphere of radius $2$.  Therefore
\begin{equation}\label{eq:hellinger-angle}
  d_F(q,q')=2\arccos\left(\sum_{k\in\B{s}}\sqrt{q(k)q'(k)}\right).
\end{equation}
For distinct point-mass limits, $d_F(\delta_k,\delta_{k'})=\pi$.  This is a
statement about the belief simplex, not about geodesics between the singular
input-output distributions $P_k$ and $P_{k'}$.

\subsection{Pythagorean identities}\label{subsec:pythagorean}

The generalized Pythagorean theorem is used in the following restricted sense.
For full-support distributions in the regular Walsh exponential family, the
standard identity
\begin{equation}\label{eq:pythagorean-interior}
  \KL(P_A\|P_C)=\KL(P_A\|P_B)+\KL(P_B\|P_C)
\end{equation}
holds whenever the e-geodesic from $A$ to $B$ is orthogonal to the m-geodesic
from $B$ to $C$.  In the CPA setting, the relevant full-support object is the
posterior belief $q_n\in Q$ as long as the prior has full support and the
Gaussian likelihood is strictly positive.  This is why projection statements in
Section~\ref{sec:cpa} are made on $Q$ or on regularized joint distributions,
not directly on disjoint-support $P_k$ endpoints.
